\documentclass[11pt,a4paper]{article}
\usepackage[margin=1.1in]{geometry}
\usepackage{amsmath,amssymb,amsthm,mathtools}
\usepackage{lmodern}
\usepackage[T1]{fontenc}
\usepackage{microtype}
\usepackage{booktabs}
\usepackage{xurl}
\input{leanlisting}
\usepackage[colorlinks,linkcolor=leankw,citecolor=leankw,urlcolor=leankw]{hyperref}
\newtheorem{theorem}{Theorem}
\newcommand{\lean}[1]{{\upshape\path{#1}}}
\newcommand{\st}[1]{{\upshape\small\textsf{[#1]}}}
\emergencystretch=2em

\title{\textbf{Fefferman Statement (B) on $\mathbb{T}^3$: A Precise Lean Formalization, a Kernel-Checked Reduction, and the Boundaries of the Frohmanian Symplectic Tether}}
\author{Benjamin Stanley Frohman\\
\small ORCID: \href{https://orcid.org/0009-0006-7068-3718}{0009-0006-7068-3718} \quad|\quad
GitHub: \href{https://github.com/BenFrohman/Clay-}{\lean{BenFrohman/Clay-}}\\
\small Concept DOI (all versions): \href{https://doi.org/10.5281/zenodo.23653012}{\texttt{10.5281/zenodo.23653012}}\\ \small Previously archived version DOI: \href{https://doi.org/10.5281/zenodo.23940667}{\texttt{10.5281/zenodo.23940667}}}
\date{Earliest Commit Provenance: 2026-10-07 06:16:45 AM CDT (UTC-5)\\
Revised edition: November 6, 2026 (America/Chicago, UTC$-6$). Not peer reviewed.}

\begin{document}
\maketitle

\begin{abstract}
This monograph records, in Lean~4 against Mathlib, a precise formal statement of Fefferman's Statement~(B) on global smoothness of the unforced incompressible Navier--Stokes equations on $\mathbb{T}^3=\mathbb{R}^3/\mathbb{Z}^3$, together with an exact ledger of the analytic inputs that the Frohmanian symplectic tether does \emph{not} supply. Statement~(B) is \emph{not} proved here. What is proved is a reduction: Statement~(B) follows from the classical local theory with the Beale--Kato--Majda continuation criterion and from the open hypothesis that the Beale--Kato--Majda integral is finite on every finite interval of existence. We also record the two-dimensional clock test: constant vortex direction ($\sin\varphi\equiv 0$) coexists with linear growth $\|\phi(t)\|_{L^\infty}=M_0^2\,t$ on steady Euler equilibria.
\end{abstract}

\section{Scope and non-claim boundary}
This repository is a boundary and gap ledger. It is \textbf{not} a resolution of the Clay Millennium Problem. Every open input is a named proposition in Lean; none is assumed as an axiom, and no proof uses \texttt{sorry}. The build and axiom audit (\lean{lean/AxiomReport.lean}, output in \lean{AXIOMS.md}) reports only \texttt{propext}, \lean{Classical.choice} and \lean{Quot.sound}. The active continuation of this programme, in which the remaining gap is decomposed into finer named statements (all open or conditional, apart from individually proved lemmas such as the planar \texttt{sinShear} obstruction; the open statement is the Frohmanian Tether Hypothesis (A)/(B), see \lean{NAMING.md}), is \href{https://github.com/BenFrohman/periodic-ns-classical-continuation}{\lean{BenFrohman/periodic-ns-classical-continuation}}.

\begin{itemize}
  \item \textbf{Repository:} \url{https://github.com/BenFrohman/Clay-}
  \item \textbf{Concept DOI:} \href{https://doi.org/10.5281/zenodo.23653012}{\texttt{10.5281/zenodo.23653012}}
  \item \textbf{Previously archived version DOI:} \href{https://doi.org/10.5281/zenodo.23940667}{\texttt{10.5281/zenodo.23940667}}
  \item \textbf{Earliest commit:} \texttt{9e61806681bfce3f0edee28293d22a6a711f0780} (2026-10-07, 06:16:45 CDT, UTC$-5$)
  \item \textbf{License:} Copyright \textcopyright\ 2026 Benjamin Stanley Frohman. Apache License, Version~2.0.
\end{itemize}

\section{Statement (B)}
Following the official problem description of Fefferman~\cite{Fef}, fix $\nu>0$ and let $u^\circ$ be a smooth, divergence-free vector field on $\mathbb{R}^3$ with $u^\circ(x+e_j)=u^\circ(x)$ for $j=1,2,3$. Take $f\equiv 0$.

\begin{quote}
\textbf{Statement (B).} There exist smooth functions $p(x,t)$, $u_i(x,t)$ on $\mathbb{R}^3\times[0,\infty)$ such that
\[
\partial_t u_i+\sum_{j=1}^3 u_j\,\partial_j u_i = \nu\,\Delta u_i-\partial_i p,\qquad \operatorname{div}u=0,\qquad u(x,0)=u^\circ(x),
\]
and $u$, $p$ are $1$-periodic in each spatial coordinate.
\end{quote}

In Lean this is the proposition \lean{Clay.StatementB.FeffermanStatementB} (file \lean{lean/Clay/StatementB.lean}). Space is $\mathbb{R}^3$ realised as Euclidean space on $\mathrm{Fin}\,3$; partial derivatives are Fr\'echet derivatives along the standard basis; the time derivative is taken within $[0,\infty)$; smoothness is $C^\infty$ on $[0,\infty)\times\mathbb{R}^3$. No component of the statement is opaque or trivial.

\section{The proved reduction}
Write $\omega=\operatorname{curl}u$. Consider the propositions
\begin{itemize}
  \item \lean{PeriodicMaximalSolutionWithBKM} (classical; Kato, Fujita--Kato, Beale--Kato--Majda~\cite{BKM}; not yet formalized): for all admissible data there are $T^\ast\in(0,\infty]$ and a smooth solution on $[0,T^\ast)$ such that $T^\ast<\infty$ implies $\int_0^{T^\ast}\|\omega(t)\|_{L^\infty}\,dt=\infty$;
  \item \lean{PeriodicBKMIntegrability} (\textbf{open}; the Clay core): every smooth solution on a finite interval $[0,T)$ satisfies $\int_0^{T}\|\omega(t)\|_{L^\infty}\,dt<\infty$.
\end{itemize}

\begin{theorem}[\lean{feffermanStatementB_of_bkm}]
\lean{PeriodicMaximalSolutionWithBKM} and \lean{PeriodicBKMIntegrability} together imply Statement~(B).
\hfill\st{Proved as an implication; Conditional on both hypotheses}
\end{theorem}
\begin{lstlisting}
theorem feffermanStatementB_of_bkm (hloc : PeriodicMaximalSolutionWithBKM)
    (hbkm : PeriodicBKMIntegrability) : FeffermanStatementB
\end{lstlisting}

\begin{theorem}[\lean{bkmIntegralFinite_of_holographic}]
If $\|\omega(t)\|_{L^\infty}$ is bounded on every finite interval of existence, then \lean{PeriodicBKMIntegrability} holds.
\hfill\st{Proved as an implication; the bound itself is open}
\end{theorem}

The second hypothesis of the first theorem is open and is equivalent in strength to Statement~(B) itself. Consequently the reduction does not prove Statement~(B).

\section{Ledger of missing inputs}
The file \lean{lean/Clay/Gaps/MissingInputs.lean} states each former placeholder as a precise proposition:
\begin{itemize}
  \item \lean{metriplecticControlsEnstrophyL2}: $\sup_{[0,T)}\int_{\mathbb{T}^3}|\omega|^2<\infty$ for every solution on a finite interval (open);
  \item \lean{czGivesBMONotLinf}: no constant $C$ gives $\|\nabla v\|_{L^\infty}\le C\|\operatorname{curl}v\|_{L^\infty}$ for smooth periodic divergence-free $v$ (classical, not formalized);
  \item \lean{spikeObstruction}: no constant $C$ gives $\|f\|_{L^\infty}\le C\|f\|_{L^6(\mathbb{T}^3)}$ for smooth periodic $f$ (classical, not formalized);
  \item \lean{holographicNotAZFCBound}: the uniform ceiling $\sup_{[0,T)}\|\omega\|_{L^\infty}<\infty$ (open);
  \item \lean{depletedStretching}: $\omega\cdot S\omega\le C|\omega|^2$ uniformly on finite intervals (open).
\end{itemize}
The former trivial definitions have been removed; the repository history retains them. For planar fields $v=(v_1(x_1,x_2),v_2(x_1,x_2),0)$ the vorticity is proved parallel to $e_3$. Kernel-checked elementary results include the identity $\Pi_u u=0$ for the pointwise tether projection, the unboundedness of the history integral $m^2t$ for $m>0$, and the conditional Riccati rearrangement.

\section{Build and axiom status}
\textbf{Proved (kernel-checked).} \texttt{lake build} of the full \texttt{Clay} library (root module and all 15 submodules, including the note \texttt{Clay.Notes.NotBlowup}) succeeds on Lean/Mathlib v4.16.0. The axiom report lists 113 theorem declarations (including compiler-generated auxiliary lemmas), none using \texttt{sorryAx}, and none using an axiom beyond \texttt{propext}, \texttt{Classical.choice}, \texttt{Quot.sound}. The substantive kernel-checked results are \lean{feffermanStatementB_of_bkm}, \lean{statementB_of_bkmIntegralFinite}, \lean{bkmIntegralFinite_of_holographic}, \lean{bkmIntegralFinite_of_vortSup_le}, the two-dimensional clock results \lean{steady_euler_clock}, \lean{alignment_does_not_delete_clock}, \lean{clock_fails_on_aligned_2d_solutions}, the planar-vorticity lemma \lean{curl_planar}, the tether identities \lean{tether_degeneracy}, \lean{tether_bracket_preservation}, and the integer rearrangement \lean{conditional_riccati}. This is a machine check of the Lean statements, not peer review.

\textbf{Conditional on.} The global-existence reduction takes \lean{PeriodicMaximalSolutionWithBKM} (classical local theory and continuation, not proved here) and \lean{PeriodicBKMIntegrability} (the still-undischarged BKM integrability assumption) as arguments. Neither is supplied by the reduction.

The separate declaration \lean{conditional_riccati} is an algebraic implication over integers. For $\mathrm{rate},C,\kappa,M\in\mathbb Z$, \lean{CorrectedMaxRate} is exactly
\[
  \mathrm{rate}\le M(CM-\kappa M^2),
\]
and the theorem rewrites its right-hand side as $CM^2-\kappa M^3$. Here \texttt{rate} is an integer variable, not a time derivative. No real-valued trajectory, differential inequality, vorticity maximum principle, or connection of this predicate to Navier--Stokes is proved. An analytical max-rate law and its PDE bridge remain separate missing inputs; this rearrangement does not supply them. Finer decompositions of the open BKM input belong to the continuation repository and are not formalized here.

\textbf{Open.} \lean{PeriodicBKMIntegrability}; every item of the ledger above marked open; and therefore Statement~(B) itself. The Clay Millennium Problem is not solved by this work.

\paragraph{An abstract criterion is not an additional analytic input.}
The new semantic regression
\lean{abstract_dhy_criterion_iff_periodicBKMIntegrability}
proves that the current \lean{dhyImpliesRegularity} proposition is equivalent to
\lean{PeriodicBKMIntegrability}: its existential packet has no connection to the
solution and can have all four functions identically zero. This is a diagnosis
of the abstract interface, not a proof of a Deng--Hou--Yu geometric criterion or
of BKM integrability. An actual geometric criterion must bind its packet to
vortex lines of the solution and justify the required quantitative estimates.

\subsection*{A solution-linked, nondegenerate vortex-line implementation}
The module \lean{Clay.Depletion.SolutionLinkedVortexLines} now supplies a
\lean{Packet u T} tied to the actual \lean{curl (u t)}. It records moving endpoints,
parameterized curves joining them, actual normalized-vorticity tangents,
nonvanishing along active segments, measurable observed endpoint magnitudes,
and an actual bound on the integral of the direction divergence. At nonzero
peak vorticity this implementation additionally requires positive segment
length: it is an explicitly nondegenerate specialization of the spatial setting
of Deng--Hou--Yu Theorem~1~\cite{DHY}, not a transcription of every degenerate
case. The paper's moving-segment Theorem~2 is for Euler and is not transferred
to viscous Navier--Stokes here.

\textbf{Proved.} The tangent has norm one wherever vorticity is nonzero. A
zero-length packet forces the actual vorticity to vanish everywhere at that
time; any nonzero actual vorticity forces positive packet length. The actual
zero solution satisfies \lean{NSOn}, admits a packet, and passes the endpoint
and BKM regressions. Given separately the pointwise spatial comparison and
finite observed endpoint integral, the new theorem proves
\[
  \Omega(t)\le K|\omega(t,x(t))|
  \le Ke^C|\omega(t,y(t))|,\qquad
  \int_0^T\Omega(t)\,dt\le Ke^C\int_0^T|\omega(t,y(t))|\,dt<\infty.
\]
This is a conditional integral consequence, not an implemented geometric
regularity criterion. Neither endpoint-integral finiteness nor BKM finiteness
is hidden in the packet fields.

\textbf{Unproved.} \lean{SpatialBridgeObligation} is a named proposition only,
with no supplied proof or added axiom. It asks for the spatial chain-rule and
fundamental-theorem argument deriving endpoint comparison from
$\operatorname{div}\omega=0$, the actual vortex-line tangent, and the integrated
$\operatorname{div}\xi$ bound. Constructing suitable packets with uniform bounds
and proving endpoint-integral finiteness for arbitrary solutions also remain
open inputs. The old free-packet equivalence remains a diagnostic, not the
interpretation of this new solution-linked data.

\section{Provenance hashes}
Source files are fingerprinted with SHA-256 (\lean{provenance/PROVENANCE_HASHES.json}) and a Merkle root (\lean{provenance/PROVENANCE_MERKLE_ROOT.json}), signed with Ed25519 key fingerprint \lean{b428fb72676f1c52}. These are file-integrity fingerprints of an earlier snapshot only; they are not a certificate of correctness or review. The authoritative correctness check is the build and axiom report above, rerun by CI on every push.

\begin{thebibliography}{9}
\bibitem{Fef} C.~L. Fefferman, \emph{Existence and smoothness of the Navier--Stokes equation}, in: The Millennium Prize Problems, Clay Mathematics Institute, 2006, pp.~57--67.
\bibitem{BKM} J.~T. Beale, T.~Kato, A.~Majda, \emph{Remarks on the breakdown of smooth solutions for the 3-D Euler equations}, Comm. Math. Phys. \textbf{94} (1984), 61--66.
\bibitem{CF} P.~Constantin, C.~Fefferman, \emph{Direction of vorticity and the problem of global regularity for the Navier--Stokes equations}, Indiana Univ. Math. J. \textbf{42} (1993), 775--789.
\bibitem{DHY} J.~Deng, T.~Y.~Hou, X.~Yu, \emph{Geometric Properties and Non-blowup of 3-D Incompressible Euler Flow}, \url{https://arxiv.org/abs/math-ph/0402032}, Theorem~1 and Lemma~1; Theorem~2 concerns Euler flow.
\end{thebibliography}

\vfill
\noindent\rule{\textwidth}{0.4pt}\\
\footnotesize
\textbf{Attribution and licensing.} Copyright \textcopyright\ 2026 Benjamin Stanley Frohman (ORCID 0009-0006-7068-3718).
Licensed under the Apache License, Version 2.0.
\end{document}
